% ===========================================================================
%  utente/08-impostazioni/00-capitolo.tex
% ===========================================================================
\chapter{Impostazioni e dati}\label{cap:impostazioni}

La voce \voce{Settings} raccoglie tre schede: \voce{Config}, \voce{Data} e
\voce{Market rules}. Sono i tre posti in cui si dice al sistema con che cosa
lavorare e con quali limiti.

\section{Data: portare dentro i prezzi}\label{sez:data}

Senza dati, tutte le pagine dell'applicazione sono vuote. È la prima cosa da
fare dopo l'installazione.

\figuraSchermata{1.0}{figure/schermate/data.png}
  {La pagina Data. A sinistra lo scarico automatico, a destra il
   caricamento a mano, sotto le categorie.}
  {data}

\subsection{Scaricare}

Si scelgono tre cose e si preme \pulsante{Download}:

\begin{description}[leftmargin = 1.4em, style = nextline, itemsep = 4pt]
  \item[\campo{Provider}] Da dove prenderli. Alcuni richiedono una chiave
    personale, che si configura dal sito.
  \item[\campo{Timeframe}] Il passo temporale. Lo stesso strumento può
    esistere a più passi: i dati giornalieri e quelli a cinque minuti
    stanno in cartelle diverse e non si danno fastidio.
  \item[\campo{Category}] Che tipo di dato è.
\end{description}

\noindent
Il lavoro parte in coda e la campanella avvisa alla fine. Sotto ai campi,
una riga dice che cosa c'è già: quanti file, quanto spazio, fino a quando
arrivano.

\subsection{Caricare un file proprio}

Il riquadro \voce{Your data} accetta CSV, Parquet e XLSX trascinati sopra.
Il sistema riconosce da solo le colonne scritte in maiuscolo, le date messe
come indice o come colonna, e i prefissi tipo \campo{adj\_}. Il riquadro
\voce{Expected format} mostra che cosa si aspetta.

\subsection{Le categorie}

Ogni categoria è una cartella, e ce ne sono di due tipi.

\begin{center}
\begin{tabular}{@{}L{22mm}L{116mm}@{}}
\toprule
\textbf{Tipo} & \textbf{Significa} \\
\midrule
\rigaTenue
\voce{1:1} & Un file per strumento, con il nome dello strumento. Il
  collegamento è automatico, e l'applicazione può dire quanti strumenti sono
  coperti e quanti no. \\
\voce{1:N} & File condivisi fra più strumenti: settori, serie
  macroeconomiche, valute. Qui non c'è nessuna copertura da calcolare, e
  deve essere la strategia a dichiarare quale file usare. \\
\bottomrule
\end{tabular}
\end{center}

\begin{nota}{«0 / 15 assets covered» non è un guasto}
Significa che la categoria esiste ma non ha ancora dati per nessuno
strumento. Si riempie scaricando, o trascinandoci dentro dei file.
\end{nota}

\section{Config: i limiti}\label{sez:config}

Sono i parametri che valgono per tutti i calcoli.

\figuraSchermata{1.0}{figure/schermate/config.png}
  {I parametri. Ogni campo porta sotto l'unità di misura, perché la
   differenza fra \num{0.12} e \num{12} è un fattore cento.}
  {config}

\begin{center}
\begin{tabular}{@{}L{44mm}L{94mm}@{}}
\toprule
\textbf{Campo} & \textbf{Che cosa limita} \\
\midrule
\multicolumn{2}{@{}l}{\textit{\small Portfolio \& allocator}}\\[2pt]
\rigaTenue
\voce{Initial capital} & Il capitale da cui partono i calcoli. \\
\voce{Annual target volatility} & Quanto si accetta che il portafoglio
  oscilli in un anno. È la manopola principale del rischio. \\
\rigaTenue
\voce{Max weight per position} & Quanto può pesare al massimo un singolo
  strumento. \\
\voce{Max gross per asset class} & Quanto può pesare al massimo una intera
  classe --- azioni, obbligazioni, valute. Serve contro una diversificazione
  solo apparente: quattro obbligazioni tutte nella stessa direzione sono una
  posizione sola. \\
\rigaTenue
\voce{Max gross leverage} & L'esposizione totale. \num{1.0} significa
  nessuna leva. \\

\addlinespace
\multicolumn{2}{@{}l}{\textit{\small Ensemble}}\\[2pt]
\rigaTenue
\voce{Target volatility} & La volatilità a cui ogni stream viene portato
  \emph{prima} di essere pesato. \\
\voce{Regime window}, \voce{thresholds} & Come si valuta lo stato del
  mercato per la modulazione. \\
\bottomrule
\end{tabular}
\end{center}

\begin{attenzione}{\voce{Target volatility} agisce anche sul conto vero}
Non è un parametro di visualizzazione. Lo stesso fattore che normalizza gli
stream nel backtest moltiplica le posizioni che il runner apre su un conto
reale. Cambiarlo cambia quanto si rischia davvero, non solo quello che si
vede sul grafico.
\end{attenzione}

\noindent
Più in basso, \voce{Asset categories} permette di aggiungere categorie di
strumenti: una volta create, una strategia può dichiararle e limitarsi a
quelle.

\section{Market rules: i regimi}\label{sez:regimi}

Un \voce{regime} è lo stato del mercato, in tre valori: \voce{bull},
\voce{neutral}, \voce{bear}. Le strategie possono dichiarare di pesare
diversamente secondo lo stato.

\figuraSchermata{1.0}{figure/schermate/regime.png}
  {I mercati configurati. Sotto, la storia degli stati: verde rialzista,
   grigio neutro, rosso ribassista.}
  {regime}

\subsection{Come si legge la pagina}

Ogni riquadro in alto è un \emph{mercato} --- un insieme di strumenti che
condividono lo stesso modo di misurare il regime:

\begin{description}[leftmargin = 1.4em, style = nextline, itemsep = 4pt]
  \item[\voce{RULE}] La regola usata per calcolare lo stato.
  \item[\voce{DRIVER}] La serie che guida la decisione --- per le azioni, lo
    SPY; per il forex, il dollaro.
  \item[\voce{SYMBOLS}] Quanti schemi di nome ricadono in questo mercato.
  \item[La barra colorata] Quanto tempo è stato passato in ciascuno stato.
\end{description}

\begin{attenzione}{«No driver available»}
Il riquadro di un mercato che non trova la propria serie guida lo dice con
un avviso, come il mercato \voce{crypto} nella schermata. Quel mercato
\textbf{non modula niente}: non è un errore che blocca qualcosa, ma le
strategie su quegli strumenti si comportano come se il regime fosse sempre
neutro. Il rimedio è scaricare la serie mancante dalla pagina \voce{Data}.
\end{attenzione}

\subsection{Il grafico}

Mostra la serie guida con lo sfondo colorato secondo lo stato. Sotto, i
conteggi: quante osservazioni, quante volte lo stato è cambiato, e la
percentuale di tempo passata in ciascuno.

\begin{nota}{Pochi cambi di stato è un bene}
Un regime deve essere uno \emph{stato}, non una reazione. Venticinque
transizioni in cinque anni vogliono dire che la regola distingue periodi
diversi; duecento vorrebbero dire che sta inseguendo il rumore, e una
strategia che cambia peso ogni settimana paga costi senza prendere
posizione.
\end{nota}
